\section*{Chapter \ref{ch:s2:credit-relative-value} --- Credit Relative Value}

\begin{solution}{exo:s2:credit-relative-value:1}
Per unit of notional, $15 - 0.9 \times 30 = -12$ basis points a year; on capital of 10\% that is $-1.2\%$ a year.
\end{solution}

\begin{solution}{exo:s2:credit-relative-value:2}
$4.5 \times 1.35\% = 6.1\%$ of the notional, $60.8\%$ of capital at ten times.
\end{solution}

\begin{solution}{exo:s2:credit-relative-value:3}
$1/0.10 = 10$ times.
\end{solution}

\begin{solution}{exo:s2:credit-relative-value:4}
On a default the protection pays par less recovery and the bond pays recovery: the position is made whole. Before that, the bond and the CDS are marked
at their spreads, and when the basis widens the bond loses more than the protection gains; with leverage, the marks can force a sale.
\end{solution}

\begin{solution}{exo:s2:credit-relative-value:5}
A high lending fee means many investors are short the bond; arbitrageurs are reluctant to hold a bond others are betting against, even hedged, so the
basis must be more negative to attract them.
\end{solution}

\begin{solution}{exo:s2:credit-relative-value:6}
A default between the two maturities: a flattener sells short-dated protection and buys long-dated protection, so a default before the short contract
matures pays both legs, which roughly offset, while the notionals, sized by spread duration, are unequal and leave a net jump-to-default exposure.
\end{solution}

\begin{solution}{exo:s2:credit-relative-value:7}
With no funding cost, selling at a cumulative loss of 50\% of capital locks in $-50.5\%$ at year 5.24, days before the basis's low; holding to the end
earned $+24.4\%$. The rule saves the book from ruin in a worse crisis and forfeits the recovery in this one; it is a leverage decision made after the
fact.
\end{solution}

\begin{solution}{exo:s2:credit-relative-value:8}
Holding to maturity requires funding to maturity. The bond is financed in short-term repo with a haircut that can rise, and margin calls on marks
force sales before maturity; the arbitrage is real only for capital that cannot be withdrawn.
\end{solution}

\begin{solution}{pb:s2:credit-relative-value:1}
\begin{enumerate}
\item CDS spread minus the bond's spread; buying the bond with protection when the bond spread is higher.
\item The hedged bond is risk-free; frictions (funding, liquidity, counterparty, collateral) keep it from being arbitraged.
\item Deviations largest for high-friction bonds, more negative with high lending fees.
\item A premium of 15 basis points, lifted to 150 in year 5 over a quarter, back over a year.
\item Index against members' CDS when the skew is wide; 55 basis points a year, Sharpe ratio 0.97.
\item Two maturities, duration-matched; 49 basis points a year, Sharpe ratio 0.75.
\item Costs, names' credit events, jump-to-default between maturities.
\item It needs forty single-name trades.
\item $+1.6\%$, $-1.1\%$ and $-3.8\%$ a year at 0, 30 and 60 basis points of funding.
\item About $-60\%$ of capital to the low; $+76\%$, $+63\%$, $+50\%$ after.
\item Leverage withdrawn by failing prime brokers; arbitrageurs forced to sell.
\item $+69.8\%$ with no funding cost, $+67.1\%$ at 30 basis points, in a year.
\item \textbf{Named result.} In normal times the negative-basis book carries 1.6\% a year on capital at zero funding cost and loses money at 30 basis
points; when funding costs jump and the basis widens from $-15$ to $-151$, it loses about 60\% of its capital at ten times, all of it marks.
\item A basis of $-15$ covers little funding.
\item Capital that was not levered to the limit before the crisis.
\item For a crisis widening at the leverage the lender will still allow then.
\item Funding and haircuts as they were, forced sales, the crisis in the sample.
\item Cross-currency credit.
\item All three are funded arbitrages whose losses come from the financing, not the hedge.
\item The cost of the balance sheet that holds the bond.
\end{enumerate}
\end{solution}

\begin{solution}{iq:s2:credit-relative-value:1}
The CDS spread minus the bond's spread on the same issuer; negative means the bond pays more than the CDS premium, so buying the bond with protection
earns a positive spread, which persists because funding and holding the bond are costly.
\end{solution}

\begin{solution}{iq:s2:credit-relative-value:2}
Buy the bond, finance it in repo, buy protection; receive the bond's coupons, pay repo and the CDS premium; on a credit event deliver the bond (or cash
settle) and receive par: the loss on the bond is covered and the position closes.
\end{solution}

\begin{solution}{iq:s2:credit-relative-value:3}
A widening of the basis with higher haircuts and repo rates, a counterparty failure on the CDS, a delivery squeeze, and the loss of the book's
financing; at ten times, the crisis widening of the chapter.
\end{solution}

\begin{solution}{iq:s2:credit-relative-value:4}
The intrinsic spread is the members' spreads weighted so that the index's premium leg matches the basket's (roughly the duration-weighted average);
the index differs because it is more liquid, trades first and is used for hedging.
\end{solution}

\begin{solution}{iq:s2:credit-relative-value:5}
Reference entities and obligations, index series and versions, recovery conventions, coupons and roll dates; at rolls the on-the-run index changes
members and maturity, and positions must be mapped to the right series.
\end{solution}

\begin{solution}{iq:s2:credit-relative-value:6}
Before default the position receives the bond spread over the risk-free rate and pays the CDS premium; at default it receives par from the protection
plus nothing from the bond beyond its recovery, which it delivers. Funding, collateral, the cheapest-to-deliver option and counterparty risk break the
equivalence.
\end{solution}
